% appendix_math.tex: the mathematics and method appendices of
% correlated_bias_prior_v2.tex.  The parent note inputs this file after
% \appendix, so it has no preamble and its \section commands are
% appendix sections.  The memory model itself (sec:memory and its
% subsections) lives in appendix_physics.tex, next to the theory
% figures; the references to it below resolve through that file.
%
% The macros are defined here only when the parent has not defined
% them already.
\providecommand{\ones}{\mathbf{1}}  % the vector of ones, (1, ..., 1)^T
\providecommand{\Prec}{\mathsf{P}}  % a projected prior precision
\providecommand{\Cov}{\mathsf{C}}   % the prior covariance of the bin biases
\providecommand{\T}{^{\mathsf{T}}}
\providecommand{\bstwo}{b_{s^2}}

\section{The accuracy test}
\label{app:bias}

This appendix gives the method of Section~\ref{sec:accuracy}: the shift
of a fit under a known drift, the two drift laws, the random truths
built on them, and the statistic $\chidof$ of Eq.~\eqref{eq:gof}, whose
value is written $\eta$.

\subsection{The shift formula}
\label{sec:biasformula}

Let $m$ be the expected data vector (the band powers), $\mathsf{C}_d$
its covariance, $\mathsf{D}$ the matrix of derivatives of $m$ with
respect to the parameters (sans-serif, not the growth factor $D$), and
$F = \mathsf{D}\T \mathsf{C}_d^{-1} \mathsf{D} + F_{\rm priors}$ the
Fisher matrix of one fit, its priors included (Appendix~\ref{sec:fisher});
under model A the parameters include the variant's higher-order
parameters, marginalized like the cosmology.  If the true biases carry
a drift the fit may not express, the data vector moves by $\delta m$,
and to first order the maximum of the posterior moves by
\begin{equation}
  \delta p \;=\; F^{-1} \mathsf{D}\T \mathsf{C}_d^{-1}\, \delta m .
  \label{eq:fisherbias}
\end{equation}
The derivation is one line: the score (the gradient of the log
posterior) at the old maximum becomes
$\mathsf{D}\T \mathsf{C}_d^{-1} \delta m$ instead of zero, and the
curvature that converts a score into a displacement is $F$.  With
$\mathsf{C}_w$ the fit's marginalized $2 \times 2$ covariance of
$(w_0, w_a)$ and $\delta w$ the $(w_0, w_a)$ entries of $\delta p$, the
shift is
\begin{equation}
  \Delta_w \;=\; \sqrt{\delta w\T\, \mathsf{C}_w^{-1}\, \delta w} ,
  \label{eq:shiftw}
\end{equation}
the displacement of the dark-energy ellipse in units of its own size;
$\Omega_m$ and $A_s$ shifts are $|\delta p|/\sigma_p$.  The threshold
drawn in the figures is $1/3$, a shift that widens the effective error
by a few percent.  Both the ellipse and the response use the same fit's
$F$, so a rigid model is punished twice: it absorbs little of the drift,
and its ellipse is small.

The drift enters through the bin biases: with $\beta_i$ the amplitudes
of Eq.~\eqref{eq:beta} and $\delta_i$ the drift's fractional deviations,
$\delta m = \sum_i (\partial m / \partial \beta_i)\, B_g\, \delta_i$.
Under model A the column $\partial m/\partial \beta_i$ is the chain-rule
column of Eq.~\eqref{eq:chain} divided by $D(\bar{z}_i)$, so a drift of
$b_1$ drags $b_2$ and $\bstwo$ along through their calibrations.  For
the growth-tied prior the response is evaluated in its own parameter
basis; for the free prior the bias derivatives span the drift exactly,
so the cosmology shift vanishes, and the accuracy of that prior is
bought entirely with the information it gives up.

\subsection{The drift laws}
\label{sec:driftlaws}

Two known drifts are injected.  The conserved-population drift applies
to every group: the deviation of Eq.~\eqref{eq:tilt}, anchored at the
group's mean redshift, nearly a straight line of a few percent across
each group.  The BGS flux-limit drift applies to BGS alone: the convex
deviation of Appendix~\ref{app:bgsflux}, a $22\%$ span whose best line
leaves a $2.4\%$ rms curvature.  Both come from
\texttt{plots/make\_theory\_data.py}; only their shapes matter, because
the free amplitudes remove the level.

\subsection{Random truths}
\label{sec:randomtruths}

A random truth sets the bias of bin $i$ in group $g$ to
\begin{equation}
  b_i \;=\; \frac{B_g}{D(\bar{z}_i)}\,
  \bigl(1 + \delta_i^{\rm drift} + \varepsilon_i\bigr)
  \label{eq:randomtruth}
\end{equation}
at the fiducial cosmology.  The drift $\delta_i^{\rm drift}$ is the
known part: the flux limit in BGS and the conserved population in the
other five groups.  The residual $\varepsilon_i$ is one draw from the
Gaussian $\mathcal{N}(0, \Sigma)$ whose covariance is the kernel of one
of the two stated priors, $\Sigma_{ij} = \sigma_g^2\,
e^{-|\bar{z}_i - \bar{z}_j|/L}$ inside a group and zero across groups.
Two families of $2000$ truths are drawn, baseline and pessimistic
scatter, and every truth is fitted under every prior, each centered on
the growth-tied shape $\delta = 0$.  No fit is told which truth it
faces: the baseline prior meets truths drawn from its own prior and
truths drawn from the pessimistic prior, and is judged on both.

Because Eq.~\eqref{eq:fisherbias} is linear in $\delta m$, and
$\delta m$ is linear in the bias deviations, the $(w_0, w_a)$ rows of
the response collapse into one $2 \times n_{\rm bins}$ matrix
$\mathsf{A}$ per prior, layout and variant, with $n_{\rm bins}$ the
number of bins of the layout:
\begin{equation}
  \delta w \;=\; \mathsf{A}\,\delta,
  \qquad
  \delta_i \;=\; \delta_i^{\rm drift} + \varepsilon_i ,
  \label{eq:linearmap}
\end{equation}
so the $2000$ truths cost one matrix product.  The draws check
themselves against the closed form
$\langle \Delta_w^2 \rangle = \Delta_{\rm drift}^2
+ \mathrm{tr}(\mathsf{C}_w^{-1} \mathsf{A} \Sigma \mathsf{A}\T)$, with
$\Delta_{\rm drift}$ the shift of the drift alone.  The response is
first order in $\delta$, while the band powers are quadratic in the
bin biases: the dropped term is a few percent of the shift for the
$1$--$3\%$ residuals and perhaps $10$--$20\%$ for the $22\%$ BGS drift.

Where the fit's prior equals the truth's scatter (the matched case),
the average shift is not a bias.  That fit's ellipse is wider than the
growth-tied one precisely because it marginalizes over the scatter the
truth now has: the fitted $(w_0, w_a)$ moves from draw to draw, and the
ellipse was built to cover those moves.  Whether an ellipse is honest
is decided by the statistic of the next subsection.

\subsection{The test statistic}
\label{sec:inflation}

A fit is correctly specified when its quoted ellipse covers the truth
with the stated probability.  The squared distance of the fitted
$(w_0, w_a)$, written $\hat{w}$, from the truth $w_{\rm true}$ in units
of the ellipse is
$\chi^2_w = (\hat{w} - w_{\rm true})\T \mathsf{C}_w^{-1}
(\hat{w} - w_{\rm true})$.  Averaged over the bias history and the data
noise,
\begin{equation}
  \langle \chi^2_w \rangle
  \;=\;
  \Delta_{\rm drift}^2
  \;+\;
  \mathrm{tr}\!\left(\mathsf{C}_w^{-1} \mathsf{A}\,\Sigma\,\mathsf{A}\T\right)
  \;+\;
  \mathrm{tr}\!\left(\mathsf{C}_w^{-1} \mathsf{C}_w^{\rm noise}\right):
  \label{eq:chi2w}
\end{equation}
the drift squared, the scatter that the fit's bias parameters, held
back by their prior, did not follow, and the noise.  Here
$\mathsf{C}_w^{\rm noise}$ is the share of $\mathsf{C}_w$ contributed by
the data noise and by the Gaussian priors other than the bias prior
(the velocity dispersions and the higher-order parameters); the cross
terms average to zero.  Eq.~\eqref{eq:chi2w} is evaluated in closed
form: it is the exact average over infinitely many universes, each with
its own noise realization and bias history.  The $2000$ draws only make
the percentile bars of Figure~\ref{fig:scatter}, and they carry no
noise.  The quoted value is $\eta = \sqrt{\langle \chi^2_w \rangle/2}$.

An honest ellipse gives $\langle \chi^2_w \rangle = 2$: if the fitted
$(w_0, w_a)$ scatters around the truth with covariance exactly
$\mathsf{C}_w$, the expected squared distance is
$\mathrm{tr}(\mathsf{C}_w^{-1}\mathsf{C}_w) = 2$, one unit per
parameter.  Growing both axes of the ellipse by $\eta$ divides
$\chi^2_w$ by $\eta^2$ and brings the average back to $2$; this matches
the coverage exactly when the real scatter of $\hat{w}$ has the shape
and orientation of the ellipse.

\paragraph{Coverage of one parameter.}
For one parameter $p$, in units of its marginalized error $\sigma_p$,
the same three terms read
\begin{equation}
  \langle \chi^2_p \rangle \;=\; \mu^2 \;+\; \mathsf{A}_p\T \Sigma\, \mathsf{A}_p \;+\; \nu ,
  \label{eq:chi2p}
\end{equation}
with $\mu$ the shift of the known drift alone, $\mathsf{A}_p$ the row
of the response that carries the bias deviations into $p$, divided by
$\sigma_p$, and $\nu$ the noise term: the $(p, p)$ entry of
$F^{-1}F_0F^{-1}$ over $\sigma_p^2$, where
$F_0 = \mathsf{D}\T\mathsf{C}_d^{-1}\mathsf{D} + F_{\rm other}$ holds the
data and the Gaussian priors other than the bias prior.  The bias prior
adds a positive semidefinite precision, $F \succeq F_0$, so
$F^{-1}F_0F^{-1} \preceq F^{-1}FF^{-1} = F^{-1}$, whose $(p,p)$ entry is
$\sigma_p^2$: $\nu \le 1$ (with $X \preceq Y$ meaning
$x\T X x \le x\T Y x$ for every $x$).  The rms displacement beyond the
noise, $(\mu^2 + \mathsf{A}_p\T\Sigma\mathsf{A}_p)^{1/2}$, is therefore
at least $(\eta^2 - 1)^{1/2}$, with $\eta = \langle \chi^2_p \rangle^{1/2}$.

Because the response is linear and both the noise and the residual are
Gaussian, the pull $\zeta = (\hat{p} - p_{\rm true})/\sigma_p$ is
exactly Gaussian, with mean $\mu$ and variance
$V = \langle \chi^2_p \rangle - \mu^2$.  The $1\sigma$ bar covers the
truth when $|\zeta| < 1$, in a fraction
\begin{equation}
  P(|\zeta| < 1) \;=\;
  \Phi\!\left(\frac{1 - \mu}{\sqrt{V}}\right)
  - \Phi\!\left(\frac{-1 - \mu}{\sqrt{V}}\right)
  \label{eq:coverage}
\end{equation}
of the universes, with
$\Phi(x) = \tfrac{1}{2}\bigl[1 + {\rm erf}(x/\sqrt{2})\bigr]$ the
cumulative distribution of a unit Gaussian.  When the prior absorbs the
known drift ($\mu \approx 0$) the coverage depends on $\eta$ alone,
\begin{equation}
  P(|\zeta| < 1) \;=\; 2\,\Phi(1/\eta) - 1 :
  \label{eq:quickcoverage}
\end{equation}
$68\%$ at $\eta = 1$, $56\%$ at $1.3$, $50\%$ at $1.5$, $38\%$ at $2$
and $31\%$ at $2.5$.  For $\eta \geq 1$, $\mu = 0$ is the most generous
split of $\langle \chi^2_p \rangle$ between drift and variance, so
otherwise Eq.~\eqref{eq:quickcoverage} is an upper bound.  For the
baseline, pessimistic and free priors ($\mu \leq 0.33$) it is exact to
$0.1\%$; for the growth-tied prior only Eq.~\eqref{eq:coverage} is
right.

Under the pessimistic scatter and model A baseline, the bar covers the
truth, on $\Delta d_{\rm com} = 62$~Mpc, $\Delta z = 0.02$, $0.05$ and
$0.2$, in $56$, $63$, $40$ and $68\%$ of universes for $\Omega_m$ and
$57$, $65$, $54$ and $68\%$ for $A_s$ under the baseline prior;
$66$--$68\%$ under the pessimistic prior; and $68.3\%$ under the free
prior.  Under the growth-tied prior it is $21$, $39$, $5$ and $7\%$ for
$\Omega_m$, where Eq.~\eqref{eq:quickcoverage} would give $24$, $39$,
$19$ and $22\%$.  For $(w_0, w_a)$ the ellipse $\chi^2_w < 2.30$
contains $68.3\%$ of a two-dimensional Gaussian, but the exact coverage
depends on how the excess of $\langle \chi^2_w \rangle$ splits between
the two directions, so only $\eta$ is quoted.

\paragraph{The baseline prior against a pessimistic truth.}
The baseline prior states that each group's residual around its free
line is small and smooth ($\sigma_g = 1\%$, $3\%$ for BGS, coherent
over $\Delta z \approx 0.2$); the pessimistic truth has $3\%$ wiggles
that change over $\Delta z \approx 0.05$.  The prior charges a large
penalty for bin biases that follow such wiggles, so the biases follow
them only in part, and the fit moves the cosmological parameters to
absorb the rest of each wiggle.  The push changes from truth to truth,
so over the $2000$ truths $\Delta_w$ spreads over the percentile bars of
Figure~\ref{fig:scatter}, and the observed universe is one of these
displacements, reported with the baseline-prior ellipse.  At $\eta = 2$,
for example, the real scatter of the best-fit $(w_0, w_a)$ is twice as
wide per axis as the reported ellipse.  The FoM, proportional to the
inverse area, would then be four times too large; the rescaled FoM is
still not a number to quote, because the ellipse's center and shape
come from a prior narrower and smoother than the truth.  If the
pessimistic truth is possible, the fit to report is the pessimistic
prior.  The opposite mismatch, a baseline truth fitted with the
pessimistic prior, gives $\eta$ below $1$: a conservative ellipse.

\paragraph{The matched case.}
The matched case lands on $1$ by an identity.  Collect all parameters
in $\theta$, with every prior centered on the fiducial
$\theta_{\rm fid}$, and let $\Delta\theta$ be the truth's departure from
it: zero in the cosmology and in every parameter with another Gaussian
prior, $B_g\,\delta_i$ in the amplitude $\beta_i$.  With $\mathbf{n}_d$
the data noise, of covariance $\mathsf{C}_d$, the linear response puts
the fit at $\hat{\theta} = \theta_{\rm fid}
+ F^{-1}\mathsf{D}\T\mathsf{C}_d^{-1}(\mathsf{D}\,\Delta\theta + \mathbf{n}_d)$.
Write $F = \mathsf{D}\T\mathsf{C}_d^{-1}\mathsf{D} + F_{\rm bias}
+ F_{\rm other}$, with $F_{\rm bias}$ the bias prior's precision (the
projection of Eq.~\eqref{eq:memproj} in the amplitude block).  Since
$F_{\rm other}\,\Delta\theta = 0$, the fit's error is
\begin{equation}
  \hat{\theta} - \theta_{\rm true}
  \;=\;
  F^{-1}\bigl(\mathsf{D}\T\mathsf{C}_d^{-1} \mathbf{n}_d
  \;-\; F_{\rm bias}\,\Delta\theta\bigr) :
  \label{eq:fiterror}
\end{equation}
the noise, plus the part of the truth's departure that the bias prior
held back.  Treating each other Gaussian prior as one more independent
measurement, the noise part has covariance $F^{-1}F_0F^{-1}$.  In the
matched case $\Prec_{\rm mem}$ is built from the same kernel $\Sigma$
as the truth's scatter (in the units of the bin amplitudes, a factor
$B_g^2$ per group), and every projection of that form obeys
\begin{equation}
  \Prec_{\rm mem}\,\Sigma\,\Prec_{\rm mem} \;=\; \Prec_{\rm mem} .
  \label{eq:matched}
\end{equation}
The proof takes one line: with $\mathsf{Q} = \Sigma^{-1}$ and
$\mathsf{M} = (\mathsf{U}\T\mathsf{Q}\mathsf{U})^{-1}$,
$\Sigma\,\Prec_{\rm mem} = \mathsf{I} - \mathsf{U}\mathsf{M}\mathsf{U}\T\mathsf{Q}$,
and $\Prec_{\rm mem}\mathsf{U} = 0$ removes the second term.  The
scatter part of Eq.~\eqref{eq:fiterror} then has covariance
$F^{-1}F_{\rm bias}F^{-1}$, the two covariances add up to
$F^{-1}FF^{-1} = F^{-1}$, and the scatter and noise terms of
Eq.~\eqref{eq:chi2w} sum to exactly $2$.  The widening the prior adds to
the ellipse equals the error the scatter inflicts.  The matched value
is $\eta = (1 + \Delta_{\rm drift}^2/2)^{1/2}$ for $(w_0, w_a)$ and
$(1 + \Delta_{{\rm drift},p}^2)^{1/2}$ for one parameter, so it departs
from $1$ only through the known drift: under model A it stays within
$3\%$ of $1$ for $(w_0, w_a)$, while a single parameter reaches $1.05$
($\Omega_m$ on $\Delta z = 0.05$, where the BGS flux limit is absorbed
least well).  Appendix~\ref{app:checks} verifies the identity
numerically.

\paragraph{When a fit is conservative.}
Write $\Prec_{\rm mem}$ in its null-space form, Eq.~\eqref{eq:memnull},
with the fit's kernel $\Sigma_{\rm fit}$, and let $\Sigma_{\rm true}$ be
the truth's.  If the truth's scatter is no larger than the fit's prior
in every combination the prior charges,
$\mathsf{N}\T\Sigma_{\rm true}\mathsf{N} \preceq
\mathsf{N}\T\Sigma_{\rm fit}\mathsf{N}$, then
$\Prec_{\rm mem}\Sigma_{\rm true}\Prec_{\rm mem} \preceq \Prec_{\rm mem}$,
the error covariance is at most $F^{-1}$, and
$\eta \le (1 + \Delta_{\rm drift}^2/2)^{1/2}$: a prior wider than the
truth wastes information but does not understate the error.  Scatter
larger or rougher than the pessimistic prior would make the
pessimistic prior's error bar too small in the same way as above; the
ladder ends at the free prior, for which $F_{\rm bias} = 0$ leaves only
the noise in Eq.~\eqref{eq:fiterror}, so $\eta = 1$ for every truth
expressible through the bin biases.

\subsection{The same test for $\Omega_m$ and $A_s$}
\label{sec:inflationomas}

The same drifts and residuals are scored for $\Omega_m$ and $A_s$, each
in units of its own marginalized error, with
$\chi^2_p = (\hat{p} - p_{\rm true})^2/\sigma_p^2$ and
$\eta = \langle \chi^2_p \rangle^{1/2}$, the case $n_{\rm dof} = 1$ of
Eq.~\eqref{eq:gof}; Eq.~\eqref{eq:coverage} converts it into coverage.
These two can react more strongly to a bias error than $(w_0, w_a)$,
because a bias error moves the clustering amplitude
$b_i^2 A_s D^2(\bar{z}_i)$ of bin $i$, which $A_s$ scales directly and
$\Omega_m$ moves through the growth factor.

\section{The projected precision for two bins}
\label{sec:twobins}

Everything essential already happens with two bins.  Call their
amplitudes $\beta_1$ and $\beta_2$, with common fiducial value $B$,
prior spread $\sigma_\beta$ per bin, and correlation coefficient $r$
between the two deviations.  The prior covariance and its inverse are
\begin{equation}
  \Cov \;=\; \sigma_\beta^2
  \begin{pmatrix} 1 & r \\ r & 1 \end{pmatrix},
  \qquad
  \Cov^{-1} \;=\; \frac{1}{\sigma_\beta^2\,(1 - r^2)}
  \begin{pmatrix} 1 & -r \\ -r & 1 \end{pmatrix}
  \label{eq:twobinC}
\end{equation}
(the off-diagonal of $\Cov\,\Cov^{-1}$ is $r - r = 0$ and the diagonal
$(1 - r^2)/(1 - r^2) = 1$).  The prior penalty of a displacement
$\Delta\beta = (\beta_1 - B,\ \beta_2 - B)\T$ is
$\Delta\beta\T\, \Cov^{-1} \Delta\beta$, and two kinds of move matter.

\paragraph{A uniform shift.}
For $\Delta\beta = c\,(1, 1)\T$, both bins up by $c$,
\begin{equation}
  \Delta\beta\T \Cov^{-1} \Delta\beta
  \;=\;
  \frac{c^2\,\bigl[(1 - r) + (1 - r)\bigr]}{\sigma_\beta^2 (1 - r^2)}
  \;=\;
  \frac{2 c^2}{\sigma_\beta^2\,(1 + r)}
  \;\neq\; 0 .
  \label{eq:shiftpenalty}
\end{equation}
The plain prior charges for a uniform shift, but a uniform shift is a
change of $B_g$, which the growth-tied model leaves free; keeping the
penalty would constrain $B_g$ and make the comparison unfair.

\paragraph{A shape change.}
For $\Delta\beta = d\,(1, -1)\T$, the bins moving apart,
\begin{equation}
  \Delta\beta\T \Cov^{-1} \Delta\beta
  \;=\;
  \frac{2 d^2}{\sigma_\beta^2\,(1 - r)} .
  \label{eq:shapepenalty}
\end{equation}
This is the penalty the prior is meant to apply: it grows without bound
as $r \to 1$, where the two bins are forced to track each other, and
relaxes to $2d^2/\sigma_\beta^2$ as $r \to 0$.  At the worked value
$r = 0.670$ of Appendix~\ref{sec:jump}, a shape change costs
$6.1\,d^2/\sigma_\beta^2$ and a uniform shift of the same size
$1.2\,c^2/\sigma_\beta^2$: the correlated prior already charges five
times more for the shape, and the projection takes the uniform charge
to zero.

\paragraph{The projected precision.}
Both rows of $\Cov^{-1}$ have the same sum, so
$\Cov^{-1}\ones = [\sigma_\beta^2 (1+r)]^{-1} (1,1)\T$ and
$\ones\T \Cov^{-1} \ones = 2 / [\sigma_\beta^2 (1+r)]$.  Substituting
into Eq.~\eqref{eq:projdef},
\begin{equation}
  \Prec \;=\;
  \frac{1}{2 \sigma_\beta^2 (1 - r)}
  \begin{pmatrix} 1 & -1 \\ -1 & 1 \end{pmatrix},
  \qquad
  \Delta\beta\T \Prec\, \Delta\beta
  \;=\;
  \frac{(\beta_1 - \beta_2)^2}{2 \sigma_\beta^2 (1 - r)} .
  \label{eq:twobinP}
\end{equation}
The projected prior charges only the difference of the two amplitudes:
a uniform shift is free, and a shape change keeps the penalty of
Eq.~\eqref{eq:shapepenalty}.  The denominator is the prior variance of
the jump, $\mathrm{Var}(\beta_1 - \beta_2) = 2\sigma_\beta^2(1 - r)$,
so the construction reduces to one sentence: \emph{the prior constrains
the jump, not the level.}

\section{The projection as a marginalization}
\label{sec:threeviews}

Equation~\eqref{eq:projdef} subtracts from $\Cov^{-1}$ a matrix built
from the single vector $\Cov^{-1}\ones$, so it removes the prior's
information about one direction, the uniform shift, and leaves
everything else alone.  Equivalently, the projected penalty of a
displacement $\Delta\beta$ is the plain penalty of what remains after
subtracting the best uniform shift, the generalized-least-squares mean
$\hat{c} = \ones\T \Cov^{-1} \Delta\beta / \ones\T \Cov^{-1} \ones$.
The derivation below shows why this is the correct prior and not just
one possible repair.  A group has $n$ bins throughout.

\subsection{Marginalizing a flat-prior amplitude}
\label{sec:viewc}

Model the group amplitudes as
\begin{equation}
  \beta \;=\; B\,\ones + \varepsilon,
  \qquad
  \varepsilon \sim \mathcal{N}(0, \Cov),
  \label{eq:genmodel}
\end{equation}
with a flat prior on the scalar $B$: the model of Eq.~\eqref{eq:model},
linearized around the fiducial.  The prior density of $\beta$ follows
by integrating $B$ out,
\begin{equation}
  p(\beta) \;\propto\;
  \int dB\,
  \exp\!\left[
    -\tfrac{1}{2}
    \bigl(\beta - B\,\ones\bigr)\T \Cov^{-1} \bigl(\beta - B\,\ones\bigr)
  \right].
  \label{eq:marg}
\end{equation}
The exponent is quadratic in $B$.  With $S = \ones\T \Cov^{-1} \ones$,
completing the square gives
\begin{equation}
  \bigl(\beta - B\ones\bigr)\T \Cov^{-1} \bigl(\beta - B\ones\bigr)
  \;=\;
  S\,\bigl(B - B_*\bigr)^2
  \;+\;
  \beta\T \Cov^{-1} \beta
  \;-\;
  \frac{\bigl(\ones\T \Cov^{-1} \beta\bigr)^2}{S},
  \qquad
  B_* = \frac{\ones\T \Cov^{-1} \beta}{S} .
  \label{eq:complsq}
\end{equation}
The $B$ integral is a Gaussian centered on $B_*$ with width
$1/\sqrt{S}$, and contributes only the $\beta$-independent constant
$\sqrt{2\pi / S}$.  What survives is
\begin{equation}
  p(\beta) \;\propto\;
  \exp\!\left[-\tfrac{1}{2}\, \beta\T \Prec\, \beta\right],
  \qquad
  \Prec \;=\; \Cov^{-1}
  - \frac{\Cov^{-1}\ones\,\ones\T \Cov^{-1}}{\ones\T \Cov^{-1} \ones},
  \label{eq:margresult}
\end{equation}
exactly the projection of Eq.~\eqref{eq:projdef}; $B_*$ is the same
generalized-least-squares mean as $\hat{c}$, because the best-fitting
amplitude dominates the integral.

The same computation has a standard matrix name.  The exponent of
Eq.~\eqref{eq:marg} is a Gaussian in the joint variables $(\beta, B)$
with the $(n+1) \times (n+1)$ precision
\begin{equation}
  \mathsf{F}_{(\beta, B)}
  \;=\;
  \begin{pmatrix}
    \Cov^{-1} & -\,\Cov^{-1}\ones \\[0.4ex]
    -\,\ones\T \Cov^{-1} & \ones\T \Cov^{-1} \ones
  \end{pmatrix}.
  \label{eq:jointprec}
\end{equation}
Removing a parameter from a precision matrix is not just deleting its
row and column: the remaining block must be corrected downward by
whatever the removed parameter could absorb.  The corrected block,
$\Prec = \mathsf{F}_{\beta\beta}
- \mathsf{F}_{\beta B}\, \mathsf{F}_{BB}^{-1}\, \mathsf{F}_{B\beta}$, is
the \emph{Schur complement} over $B$, and it reproduces
Eq.~\eqref{eq:projdef} term by term.  So $\Prec$ is precisely the
statement ``\emph{$B_g$ is sampled with a flat prior and
marginalized}'', folded into one matrix.

\section{Many bins: the tridiagonal inverse}
\label{sec:general}

\subsection{Three bins: total correlation versus direct links}
\label{sec:kernel}

A group has $n$ bins at redshifts $\bar{z}_1 < \dots < \bar{z}_n$,
unevenly spaced in general, with the prior covariance
$\Cov_{ij} = \sigma_\beta^2 \exp(-|\bar{z}_i - \bar{z}_j|/L)$ of
Eq.~\eqref{eq:covdef}.  This matrix is positive definite for distinct
redshifts, so $\Cov^{-1}$ exists.  Its structure shows on three evenly
spaced bins: with neighbor correlation $r = e^{-\Delta z / L}$ the outer
pair carries $r^2$, and
\begin{equation}
  \Cov \;=\; \sigma_\beta^2
  \begin{pmatrix} 1 & r & r^2 \\ r & 1 & r \\ r^2 & r & 1 \end{pmatrix},
  \qquad
  \Cov^{-1} \;=\; \frac{1}{\sigma_\beta^2\,(1 - r^2)}
  \begin{pmatrix} 1 & -r & 0 \\ -r & 1 + r^2 & -r \\ 0 & -r & 1
  \end{pmatrix}.
  \label{eq:threebin}
\end{equation}
Row 1 of $\Cov^{-1}$ times column 3 of $\Cov$ gives, without the
prefactors, $1 \cdot r^2 - r \cdot r + 0 \cdot 1 = 0$, as the
off-diagonal of the identity must.  Bins 1 and 3 \emph{are} correlated
($\Cov$ contains $r^2$), yet their entry in $\Cov^{-1}$ is zero: $\Cov$
gives the total correlation, what $\delta_1$ says about $\delta_3$ when
nothing else is known, while $\Cov^{-1}$ gives the direct one, what
$\delta_1$ still says once $\delta_2$ is known, and that is nothing.
The exponential kernel makes this exact because its correlations
multiply along the chain,
\begin{equation}
  e^{-(\bar{z}_3 - \bar{z}_1)/L}
  \;=\;
  e^{-(\bar{z}_3 - \bar{z}_2)/L}\; e^{-(\bar{z}_2 - \bar{z}_1)/L} ,
  \label{eq:factorize}
\end{equation}
so the middle bin already carries everything the outer bins share.  A
Gaussian kernel $e^{-\Delta z^2 / L^2}$ does not factorize this way,
because squared separations do not add along the chain, and its
inverse is dense.

\subsection{The chain of links}
\label{app:chain}

Because correlations multiply along the chain, the prior can be built
one bin at a time: the joint density of the deviations factorizes into
neighbor conditionals,
\begin{equation}
  \begin{gathered}
    p(\delta_1, \dots, \delta_n)
    \;=\;
    p(\delta_1)\, p(\delta_2 \mid \delta_1) \cdots
    p(\delta_n \mid \delta_{n-1}),
    \\
    \bigl\langle \delta_{k+1} \mid \delta_k \bigr\rangle
    = r_k\,\delta_k,
    \qquad
    \mathrm{Var}\bigl(\delta_{k+1} \mid \delta_k\bigr)
    = \sigma_g^2\,\bigl(1 - r_k^2\bigr),
  \end{gathered}
  \label{eq:chainfactor}
\end{equation}
with $r_k = e^{-(\bar{z}_{k+1} - \bar{z}_k)/L}$ the correlation of the
neighbors $k$ and $k+1$.  The conditional follows from the covariances
alone.  Split $\delta_{k+1} = \lambda\,\delta_k + e$ and choose $\lambda$
so that the leftover $e$ carries no correlation with $\delta_k$:
$\mathrm{Cov}(e, \delta_k) = \sigma_g^2 r_k - \lambda\,\sigma_g^2 = 0$
gives $\lambda = r_k$, and then
$\mathrm{Var}(e) = \sigma_g^2 - 2 r_k^2 \sigma_g^2 + r_k^2 \sigma_g^2
= \sigma_g^2 (1 - r_k^2)$.  For jointly Gaussian variables uncorrelated
means independent, so knowing the neighbor explains the fraction
$r_k^2$ of the bin's variance and the fraction $1 - r_k^2$ is new.
Earlier bins add nothing: for $j \le k$,
\begin{equation}
  \mathrm{Cov}(e, \delta_j)
  \;=\;
  \sigma_g^2\, e^{-(\bar{z}_k - \bar{z}_j)/L}
  \left[
    e^{-(\bar{z}_{k+1} - \bar{z}_k)/L} - r_k
  \right]
  \;=\; 0 ,
  \label{eq:pastfree}
\end{equation}
because the correlation factorizes through $\bar{z}_k$.  Taking
$-2 \ln p$ gives the penalty as a sum over neighbor links,
\begin{equation}
  \delta\T \bigl(\sigma_g^2 \mathsf{R}\bigr)^{-1} \delta
  \;=\;
  \frac{\delta_1^2}{\sigma_g^2}
  \;+\;
  \sum_{k=1}^{n-1}
  \frac{\bigl(\delta_{k+1} - r_k\,\delta_k\bigr)^2}
       {\sigma_g^2\,\bigl(1 - r_k^2\bigr)} :
  \label{eq:chainsum}
\end{equation}
each term compares one bin with $r_k$ times the bin before it.
\emph{Directly}, only neighbors talk; \emph{in total}, the influence
propagates down the whole chain.  Probability theory calls this the
Markov property.

\subsection{The tridiagonal inverse}
\label{sec:tridiag}

Expanding each square of Eq.~\eqref{eq:chainsum} and collecting the
coefficient of every product $\delta_i\delta_j$ reads off the precision
matrix.  $\delta_1^2$ appears in the first term and in the $k = 1$
link as the predictor, $1 + r_1^2/(1 - r_1^2) = 1/(1 - r_1^2)$; an
interior $\delta_k^2$ appears in link $k-1$ as the predicted bin and in
link $k$ as the predictor; the last bin is only ever predicted.  Each
product $\delta_k\delta_{k+1}$ appears in its own link only, and no
product of non-neighbors appears.  For the amplitudes
($\Delta\beta = B_g\,\delta$, $\sigma_\beta = \sigma_g B_g$) the inverse
is therefore tridiagonal,
\begin{equation}
  \sigma_\beta^2\,\bigl(\Cov^{-1}\bigr)_{kk} =
  \begin{cases}
    \dfrac{1}{1 - r_1^2}, & k = 1,\\[2.2ex]
    \dfrac{1}{1 - r_{k-1}^2} + \dfrac{r_k^2}{1 - r_k^2}, & 1 < k < n,\\[2.2ex]
    \dfrac{1}{1 - r_{n-1}^2}, & k = n,
  \end{cases}
  \qquad
  \sigma_\beta^2\,\bigl(\Cov^{-1}\bigr)_{k, k+1} = \frac{-\,r_k}{1 - r_k^2}.
  \label{eq:uneven}
\end{equation}
The off-diagonal entry is the link of pair $k$, stronger when the bins
are closer; the diagonal entry of a bin collects the links touching
it.  For even spacing, all $r_k = r$, the interior diagonal becomes
$(1 + r^2)/(1 - r^2)$ and
\begin{equation}
  \Cov^{-1} \;=\;
  \frac{1}{\sigma_\beta^2\,(1 - r^2)}
  \begin{pmatrix}
    1      & -r     &        &        &        \\
    -r     & 1+r^2  & -r     &        &        \\
           & -r     & \ddots & \ddots &        \\
           &        & \ddots & 1+r^2  & -r     \\
           &        &        & -r     & 1
  \end{pmatrix},
  \label{eq:ar1}
\end{equation}
which for $n = 3$ is Eq.~\eqref{eq:threebin}.

\subsection{The projection, group by group, and its limits}
\label{sec:projgroups}

Each group receives the projected precision of Eq.~\eqref{eq:projdef},
built from its own $\Cov$; across groups the prior is zero, so the
total prior precision is block diagonal.  $\Prec\,\ones = 0$, $\Prec$ is
positive semidefinite (no displacement is rewarded), and it has rank
$n - 1$: of the $n$ directions a group can move in, exactly the uniform
shift is free.

\label{sec:limits}
The kernel model contains both extremes.  As $L \to \infty$, every
$r_k \to 1$, the jump variance $2\sigma_g^2(1 - r_k)$ of every
neighbor pair vanishes, and only the uniform shift is left: the
growth-tied model, whatever $\sigma_g$.  As $\sigma_g \to \infty$ at
fixed $L$, $\Cov^{-1} \to 0$ and $\Prec \to 0$: the free model.

\section{The memory projection, step by step}
\label{app:memory}

Throughout, a group has $n$ bins, $\Cov$ is the kernel covariance
$\Cov_{ij} = \sigma_g^2 e^{-|\bar{z}_i - \bar{z}_j|/L}$ of the
residuals, $\mathsf{Q} = \Cov^{-1}$, $\mathsf{U} = [\,\ones,\ v\,]$ is the
$n \times 2$ design matrix of Eq.~\eqref{eq:memproj} with $v$ the
centered redshifts, and $\theta = (c, s_g)\T$ collects a line's level
and slope.  (For the bin amplitudes every covariance carries a factor
$B_g^2$ that cancels in each identity below.)

\subsection{The line fit, the split, and the marginalization}

Minimizing the penalty of a displacement $\delta$ over the line,
$\chi^2(\theta) = (\delta - \mathsf{U}\theta)\T \mathsf{Q}\,
(\delta - \mathsf{U}\theta)$, gives the normal equations of a
generalized least-squares fit,
\begin{equation}
  \mathsf{U}\T \mathsf{Q}\,\mathsf{U}\;\hat{\theta}
  \;=\;
  \mathsf{U}\T \mathsf{Q}\,\delta ,
  \qquad
  \hat{\theta} = (\hat{c}, \hat{s}_g)\T ,
  \label{eq:memnormal}
\end{equation}
with $\mathsf{U}\T\mathsf{Q}\mathsf{U}$ positive definite whenever the
group has two distinct redshifts.  Writing
$\delta = \mathsf{U}\hat{\theta} + \delta_\perp$, the normal equations
say $\mathsf{U}\T \mathsf{Q}\, \delta_\perp = 0$: the residual is
orthogonal to both columns in the $\mathsf{Q}$ metric.  The cross terms
vanish and the penalty splits,
\begin{equation}
  \delta\T \mathsf{Q}\,\delta
  \;=\;
  \hat{\theta}\T
  \bigl(\mathsf{U}\T \mathsf{Q}\,\mathsf{U}\bigr)\hat{\theta}
  \;+\;
  \delta_\perp\T \mathsf{Q}\, \delta_\perp ,
  \qquad
  \delta\T \Prec_{\rm mem}\,\delta
  \;=\;
  \delta_\perp\T \mathsf{Q}\, \delta_\perp :
  \label{eq:memsplit}
\end{equation}
the projection keeps the residual's penalty and discards the line's.
A pure line has $\delta_\perp = 0$, so $\Prec_{\rm mem}\mathsf{U} = 0$;
the residuals obey two linear constraints, so $\Prec_{\rm mem}$ has rank
$n - 2$; and $\delta_\perp\T \mathsf{Q}\, \delta_\perp \geq 0$ makes it
positive semidefinite.

Marginalization gives the same matrix, as for the single column
$\ones$ of Appendix~\ref{sec:threeviews}.  With
$\delta \sim \mathcal{N}(\mathsf{U}\theta, \Cov)$ and
$\theta \sim \mathcal{N}(0, \Lambda)$, $\Lambda = \mathrm{diag}(W^2,
\tau^2)$ the prior variances of the level and the slope, the marginal
covariance of $\delta$ is $\Cov + \mathsf{U}\Lambda\mathsf{U}\T$, and the
Woodbury identity inverts it:
\begin{equation}
  \bigl(\Cov + \mathsf{U}\Lambda\mathsf{U}\T\bigr)^{-1}
  \;=\;
  \mathsf{Q} - \mathsf{Q}\,\mathsf{U}
  \bigl(\Lambda^{-1} + \mathsf{U}\T \mathsf{Q}\,\mathsf{U}\bigr)^{-1}
  \mathsf{U}\T \mathsf{Q} .
  \label{eq:woodbury}
\end{equation}
As the prior widths grow, $\Lambda^{-1} \to 0$ and the right-hand side
becomes $\Prec_{\rm mem}$; the leading correction carries
$\Lambda^{-1}$, so the deviation falls as $1/W^2$.

\subsection{The null-space form}

Let $\mathsf{N}$ be an $n \times (n-2)$ matrix whose columns span the
combinations of bins that no straight line can produce:
$\mathsf{N}\T \mathsf{U} = 0$, plain Euclidean orthogonality to $\ones$
and to $v$.  Then
\begin{equation}
  \mathsf{Q} - \mathsf{Q}\,\mathsf{U}
  \bigl(\mathsf{U}\T \mathsf{Q}\,\mathsf{U}\bigr)^{-1}
  \mathsf{U}\T \mathsf{Q}
  \;=\;
  \mathsf{N}
  \bigl(\mathsf{N}\T \Cov\, \mathsf{N}\bigr)^{-1}
  \mathsf{N}\T .
  \label{eq:memnull}
\end{equation}
Both sides annihilate $\mathsf{U}$: the left by construction, the right
because $\mathsf{N}\T\mathsf{U} = 0$.  Both sides map $\Cov\mathsf{N}$
to $\mathsf{N}$: the right because
$\mathsf{N}(\mathsf{N}\T\Cov\mathsf{N})^{-1}
(\mathsf{N}\T\Cov\mathsf{N}) = \mathsf{N}$, the left because
$\mathsf{Q}\Cov\mathsf{N} = \mathsf{N}$ and the subtracted term meets
$\mathsf{U}\T\mathsf{N} = 0$.  And the columns of $\Cov\mathsf{N}$ and
of $\mathsf{U}$ together span the whole space: if
$\Cov\mathsf{N}x + \mathsf{U}y = 0$, multiplying by $\mathsf{N}\T$ gives
$\mathsf{N}\T\Cov\mathsf{N}\,x = 0$, so $x = 0$, and then $y = 0$.  Two
matrices that agree on a spanning set are equal.  The prior therefore
reads off the combinations $\mathsf{N}\T\delta$ that no line can make
and charges each by the inverse of its own prior covariance.  With the
single column $\mathsf{U} = \ones$ the same identity gives the kernel
projection of Eq.~\eqref{eq:projdef}, with $\mathsf{N}$ spanning the
differences between bins.

\subsection{Worked examples}

One or two bins are fully free.  For $n = 2$ the matrix $\mathsf{U}$ is
square and invertible, so the subtracted term collapses to $\mathsf{Q}$
and the projection vanishes; for $n = 1$, $\mathsf{N}$ has no columns.
Three evenly spaced bins show what survives.  With spacing $\Delta$ and
$r = e^{-\Delta/L}$, $v = (-\Delta, 0, \Delta)\T$, and the only direction
orthogonal to $\ones$ and $v$ is the discrete curvature
$\mathsf{N} = (1, -2, 1)\T$, whose prior variance is
$\mathsf{N}\T \Cov\, \mathsf{N} = \sigma_g^2(6 - 8r + 2r^2) =
2\sigma_g^2(1 - r)(3 - r)$ (the diagonal gives $1 + 4 + 1$, the two
neighbor pairs $-8r$, the outer pair $+2r^2$).  The penalty is
\begin{equation}
  \delta\T \Prec_{\rm mem}\,\delta
  \;=\;
  \frac{\bigl(\delta_1 - 2\delta_2 + \delta_3\bigr)^2}
       {2\,\sigma_g^2\,(1 - r)(3 - r)} :
  \label{eq:memcurv}
\end{equation}
the memory model charges only the curvature and is blind to any
straight-line drift, where the kernel model of
Eq.~\eqref{eq:threebin} charges the two neighbor differences and so
taxes a drift as heavily as a zigzag of the same step.

\subsection{The universal-kriging predictor}

Predict bin $k+1$ from the first $k \geq 2$ bins, with flat priors on
$\theta$.  Write $\Cov_k$ for the past kernel block, $\mathsf{U}_k$ for
the past rows of $\mathsf{U}$, $u_{k+1} = (1, v_{k+1})\T$ for the next
row, and $\gamma$ for the vector of kernel covariances between bin $k+1$
and the past.  Given $\theta$, the Gaussian conditional is
\begin{equation}
  \mathrm{E}\bigl[\delta_{k+1} \mid \delta_{1..k}, \theta\bigr]
  \;=\;
  u_{k+1}\T\theta
  \;+\;
  \gamma\T \Cov_k^{-1}
  \bigl(\delta_{1..k} - \mathsf{U}_k\,\theta\bigr).
  \label{eq:memcond}
\end{equation}
With a flat prior the posterior of $\theta$ is centered on the fit
$\hat{\theta}$ of Eq.~\eqref{eq:memnormal} applied to the past block,
and Eq.~\eqref{eq:memcond} is linear in $\theta$, so averaging replaces
$\theta$ by $\hat{\theta}$.  For the exponential kernel the weights
collapse: by Eq.~\eqref{eq:factorize} every entry of $\gamma$ is $r_k$
times the corresponding entry of $\Cov_k$'s last column, so
$\gamma\T \Cov_k^{-1} = r_k\, e_k\T$, with $e_k$ the $k$-th unit vector.
Substituting gives Eq.~\eqref{eq:memkriging}: the fitted line at the
new bin plus one Markov step from the last residual.

\subsection{The finite-slope correlation}

With a Gaussian slope of variance $\tau^2$ and the level held in $B_g$,
the deviations have the covariance of Eq.~\eqref{eq:memcov}, so
$\mathrm{Var}(\delta_i) = \tau^2 v_i^2 + \sigma_g^2$.  For separations
much larger than $L$ the kernel term is exponentially small, and
dropping it gives the correlation of Eq.~\eqref{eq:memrho}.  When
$\tau |v_i|$ and $\tau |v_j|$ both exceed $\sigma_g$, the $\sigma_g^2$
terms become negligible and the coefficient tends to
$v_i v_j / |v_i v_j| = \pm 1$.

\section{How the prior enters the Fisher matrix}
\label{sec:fisher}

A Gaussian prior adds its precision to the Fisher matrix,
\begin{equation}
  F \;=\; F_{\rm data} \;+\; F_{\rm priors},
  \label{eq:fisher}
\end{equation}
with $F_{\rm priors}$ holding one block per prior.  The projected
precision of each group, $\Prec_{\rm mem}$ of Eq.~\eqref{eq:memproj}
for the baseline and pessimistic priors, is added to the rows and
columns of that group's amplitudes $\beta_i$; all other parameters are
unchanged.  The growth-tied prior is not a block of this kind: it sums
each group's amplitude columns into one (Appendix~\ref{sec:ladderimpl}).

Bins that lost all their band powers to the scale cut have zero rows in
$F_{\rm data}$.  In the free model such a parameter is unconstrained and
must be dropped before inversion.  With the correlated prior it stays:
the prior ties the empty bin to its neighbors, so the matrix stays
invertible as long as the group's data fix its free level and tilt.
Appendix~\ref{app:checks} describes the pseudo-inverse that guards the
one system where a direction carries no information.

\section{Numerical verification}
\label{app:verify}

The companion script \texttt{verify\_prior.py}, in this folder, checks
every identity of these appendices with the cocoa Python, $23$ checks
that each print one PASS or FAIL line with the measured deviation: the
two-bin, three-bin, even and uneven tridiagonal inverses and the
projection's null direction, semidefiniteness and rank; the
Pythagorean split and the Schur complement; the neighbor-prediction
leftover and, by Monte Carlo, the jump variance with the worked example
of Appendix~\ref{sec:jump}; the continuity equation against
Eq.~\eqref{eq:fry}, the conserved deviations and the neighbor rms
values; and for the memory model the annihilation of $\mathsf{U}$ with
rank $n - 2$, the $1/W^2$ Woodbury limit, the null-space identity, the
three-bin curvature penalty, the vanishing for one and two bins, the
kriging predictor, the end-to-end correlation, the straight-line fits
of Appendix~\ref{sec:memphysics} and the LRG3 magnitude-limit rate.  The
theory figures are drawn from data written by
\texttt{plots/make\_theory\_data.py}, from the same flat $\Lambda$CDM
growth factor and fiducial amplitudes.
